% =====================================================================
% CAPITOLO 4 — IMPLEMENTAZIONE DELLA SUITE DI AGENTI
% =====================================================================
% POLPA, parte A2 di B.
% Scopo: dettagli tecnici sufficienti a convincere della SOLIDITA'
% delle realizzazioni, senza diventare documentazione operativa.
% =====================================================================

\chapter{La suite di agenti: implementazione}
\label{chap:agenti-implementazione}

Mentre il Capitolo~\ref{chap:agenti-architettura} ha presentato i principi guida
e le scelte progettuali ad alto livello, questo capitolo descrive il modo in cui
tali scelte sono state effettivamente realizzate. L'approccio è diviso in due fasi: 
dapprima si descrivono gli elementi trasversali alla suite (stack
tecnologico, organizzazione del codice, infrastruttura come codice e catena di
sicurezza); successivamente, ciascuno dei quattro agenti viene esaminato nei
suoi dettagli implementativi rilevanti, con particolare attenzione alle decisioni
non banali e a quelle che hanno richiesto un adattamento rispetto alla
documentazione standard delle librerie e dei servizi cloud impiegati. Il livello
di dettaglio scelto è quello sufficiente a giustificare la solidità delle
realizzazioni, non quello di una documentazione operativa.

% =====================================================================
\section{Stack tecnologico, infrastruttura, sicurezza}
\label{sec:agenti-stack}
% =====================================================================

\subsection{Stack tecnologico}
\label{subsec:agenti-stack-tech}

L'implementazione della suite di agenti è basata su uno stack cloud-native
omogeneo, scelto per minimizzare la varianza operativa fra servizi e per
delegare ai \emph{managed services} di Google Cloud Platform tutti i compiti
infrastrutturali a basso valore aggiunto.

\begin{itemize}
  \item \textbf{Backend.} Tutti gli agenti sono scritti in Python~3.10+ ed
        esposti tramite il framework asincrono FastAPI. La scelta è motivata
        dalla natura prevalentemente \emph{I/O bound} delle interazioni con i
        Large Language Model: la generazione di una risposta richiede da
        centinaia di millisecondi a decine di secondi, quasi interamente spesi
        in attesa del modello, dei motori di retrieval o di BigQuery. Il modello
        di concorrenza \texttt{async/await} consente di sostenere centinaia di
        richieste concorrenti per istanza Cloud Run senza ricorrere a
        thread o processi multipli. L'accesso al database relazionale avviene
        in modo asincrono tramite SQLAlchemy~2 e il driver
        \texttt{asyncpg}, con \emph{pool} pre-pingati e ricicli automatici per
        sopravvivere alle finestre di manutenzione del DBaaS.
  \item \textbf{Frontend.} Le interfacce sono React~19 con TypeScript,
        impacchettate con Vite e stilizzate con Tailwind~CSS. La scelta
        di un \emph{bundler} ESM moderno è funzionale alla pipeline CI/CD:
        il build è di pochi secondi anche su immagini di partenza minimali,
        riducendo i tempi di rilascio.
  \item \textbf{Database e Data Warehouse.} Cloud SQL per
        PostgreSQL~15~\cite{gcp_cloudsql} è il \emph{system of record}
        applicativo: vi risiedono utenti, conversazioni, messaggi, quote,
        ACL, audit log e parametri di configurazione. Google
        BigQuery~\cite{gcp_bigquery} è invece il motore analitico
        utilizzato dagli agenti Text-to-SQL (AllyCare e SysAid), con
        dataset alimentati in modo continuo da flussi appositi.
  \item \textbf{Large Language Models.} L'implementazione poggia su due
        famiglie di modelli erogati come \emph{managed APIs}: Gemini~2.5
        Flash e Pro su Vertex AI~\cite{gcp_vertexai}, e Claude Sonnet~4.6
        attraverso l'SDK Anthropic veicolato dal canale
        \texttt{AnthropicVertex}~\cite{anthropic_docs}. La duplicazione
        deliberata in SysAid (Sez.~\ref{sec:agenti-sysaid-impl}) è la
        prima leva concreta del principio di agnosticità descritto nel
        Capitolo~\ref{chap:agenti-architettura}.
  \item \textbf{Orchestrazione di agenti.} Viene utilizzato il \emph{Google Agent
        Development Kit}~(ADK)~\cite{google_adk_docs}, una libreria
        ufficiale che astrae la gestione della sessione, l'invocazione
        ciclica del modello e l'integrazione di \emph{tool} basati su
        \emph{function calling}. 
 \item \textbf{Infrastruttura.} Ogni componente è un servizio
        serverless eseguito su Google Cloud Run~\cite{gcp_cloudrun},
        descritto interamente da codice Terraform~\cite{terraform_docs}.
\end{itemize}

\subsection{Organizzazione del codice}
\label{subsec:agenti-monorepo}

Il codice sorgente dell'intera suite vive in un unico monorepo Bitbucket,
suddiviso in directory autoconsistenti, una per servizio. Ogni directory
contiene il proprio \texttt{Dockerfile}, le proprie dipendenze e il proprio
file di descrizione della pipeline (\texttt{cloudbuild.yaml}).

\begin{figure}[h]
  \centering
  \begin{lstlisting}[language=bash]
c30-ai-agents/
  vera-backend/         # Vera (FastAPI + ADK + Gemini)
  vera-frontend/
  docs-backend/         # CooPolicy (FastAPI + Vertex AI RAG Engine)
  nps-backend/          # AllyCare (FastAPI + ADK + Gemini + BigQuery)
  nps-frontend/
  sysaid-backend/       # SysAid Gemini (FastAPI + Vertex AI)
  sysaid-sonnet-backend/# SysAid Claude (FastAPI + Anthropic SDK)
  sysaid-frontend/
  portal-frontend/      # Portale unico di accesso (MSAL)
  admin-backend/        # Pannello amministrativo
  admin-frontend/
  iac-infrastructure/   # Terraform + Cloud Build
  ai-dev-documents/     # Documentazione architetturale
  \end{lstlisting}
  \caption[Struttura del monorepo]{Struttura del monorepo. Ogni servizio è una
  unit\`a di deploy autonoma; \texttt{iac-infrastructure} descrive l'ambiente
  GCP nella sua interezza.}
  \label{fig:monorepo_tree}
\end{figure}

I backend più maturi (Vera e Admin) adottano una separazione interna per
\emph{concerns}: la directory \texttt{api/} contiene esclusivamente i router
FastAPI; \texttt{services/} le funzioni di dominio prive di stato HTTP;
\texttt{models/} le entit\`a SQLAlchemy; \texttt{schemas/} i modelli Pydantic
per validazione di input e output; \texttt{core/} le utility trasversali
(database, dipendenze, sicurezza). I backend più piccoli (CooPolicy, AllyCare e
le due varianti SysAid) condividono la stessa filosofia in forma compressa, con
un singolo \texttt{main.py} che monta i \emph{router} e una manciata di moduli
ausiliari (\texttt{auth.py}, \texttt{tools.py}, \texttt{agent.py}). Questa
deliberata asimmetria è stata mantenuta perché il \emph{churn} di codice in
quegli agenti è basso: introdurre prematuramente l'astrazione a quattro
directory sarebbe stato un costo senza beneficio.

\subsection{Infrastruttura come codice e pipeline CI/CD}
\label{subsec:agenti-iac}

L'intera infrastruttura GCP è descritta in Terraform~\cite{terraform_docs} ed è
parametrizzata su due ambienti, \texttt{test} e \texttt{prod}, gestiti come
\emph{workspace} indipendenti con \emph{state} remoto su un bucket Cloud
Storage. Le risorse provisionate comprendono: l'abilitazione delle API GCP
necessarie, le istanze Cloud Run con i loro service account dedicati, l'istanza
Cloud SQL con il relativo \emph{Private Service Connect} consumer endpoint, le
politiche di Identity and Access Management (IAM), il bucket di Artifact
Registry per le immagini Docker, il template Model Armor, la rete (subnet,
indirizzi IP riservati, regole di forwarding) e il bilanciatore di carico HTTPS
interno con il certificato TLS recuperato da Secret Manager.

Le decisioni di sicurezza fondamentali sono codificate direttamente nel
modulo Terraform e dunque non sono modificabili senza una \emph{pull request}
revisionata:

\begin{enumerate}
  \item \textbf{Egress limitato a indirizzi privati.} Ogni servizio Cloud Run
        è collegato alla VPC interna mediante \emph{Direct VPC Egress} con
        \texttt{egress = PRIVATE\_RANGES\_ONLY}: il container può raggiungere
        gli endpoint privati GCP (Vertex AI, BigQuery, Cloud SQL, RAG Engine)
        ma non internet pubblico. Il vincolo è particolarmente importante per
        gli agenti Text-to-SQL, dove un'eventuale estrazione di dati verso
        l'esterno via DNS o webhook è di fatto resa impossibile a livello di
        rete.
  \item \textbf{Database non esposto su IP pubblico.} L'istanza Cloud SQL ha
        \texttt{ipv4\_enabled = false}: la connessione passa solo attraverso
        \emph{Private Service Connect} (PSC), con un \emph{forwarding rule}
        regionale verso il \emph{service attachment} di Cloud SQL.
  \item \textbf{Bilanciatore di carico interno.} Il \emph{frontend} è
        accessibile solo via un \emph{Application Load Balancer regionale
        interno} (\texttt{INTERNAL\_MANAGED}) che termina TLS con un
        certificato wildcard custodito in Secret Manager. L'esposizione
        avviene perciò solo all'interno della rete aziendale o, per gli
        utenti remoti, dietro VPN, soddisfacendo il requisito di
        accesso privato (Cap.~\ref{chap:agenti-architettura},
        Sez.~\ref{sec:agenti-principi}).
  \item \textbf{Segregazione di credenziali.} Tutte le password generate
        (database, JWT) sono prodotte da \texttt{random\_password},
        depositate in Secret Manager e iniettate nei container via
        \texttt{value\_source.secret\_key\_ref}: nessun segreto
        attraversa mai il file system dello sviluppatore.
\end{enumerate}

La pipeline CI/CD si attiva ad ogni \emph{push}: in ambiente di test sul branch
\texttt{test}, in produzione sul branch \texttt{main}. Cloud Build, configurato
tramite un \emph{trigger} per ciascun servizio collegato al rispettivo
repository Bitbucket, esegue i passi descritti nel \texttt{cloudbuild.yaml}
locale: build dell'immagine Docker, push su Artifact Registry e aggiornamento
dell'immagine attiva del relativo servizio Cloud Run. Per l'ambiente di
produzione, il \emph{trigger} ha \texttt{approval\_required = true}, che
implementa lo step approvativo richiesto dal principio di
\emph{Automazione dei rilasci} (Cap.~\ref{chap:agenti-architettura}).

\subsection{Identità, autenticazione e autorizzazione}
\label{subsec:agenti-auth}

L'autenticazione è interamente delegata a Microsoft Entra ID, che è il
\emph{tenant} identitario già adottato dalla cooperativa. La scelta del flusso
e della libreria segue lo standard di mercato per le \emph{single-page
applications}: il browser ottiene un \emph{access token} via
OAuth~2.0~\cite{rfc6749oauth2} e OpenID Connect~\cite{oidc1_0core} con
\emph{Authorization Code Flow + PKCE}, gestito dalla
\emph{Microsoft Authentication Library} (MSAL)~\cite{msft_msal} caricata nel
portale e nel pannello di amministrazione. Il token è un JWT firmato in
RS256~\cite{rfc7519jwt} ed è il vettore di autorizzazione che attraversa la
suite.

Sul lato server, ogni backend monta una dipendenza FastAPI \texttt{verify\_token}
che esegue le verifiche minime ma non negoziabili: scarica e mette in cache
le chiavi pubbliche da \texttt{JWKS} esposto da Microsoft, identifica la
chiave corretta tramite il campo \texttt{kid} dell'header, decodifica il
token con la libreria \texttt{python-jose} validandone audience e
firma e, infine, valida manualmente l'\texttt{issuer} contro la coppia
\texttt{(sts.windows.net, login.microsoftonline.com)}: questa
ridondanza è necessaria perché Entra ID emette token con \texttt{iss}
v1 o v2 a seconda dell'\emph{endpoint} usato, e l'opzione \texttt{verify\_iss}
del decoder accetta un solo valore alla volta. Una rotazione di chiavi non
prevista forza un secondo \emph{fetch} del JWKS prima di restituire un
errore 401: la cache non è quindi un \emph{single point of failure}.

L'autorizzazione fine-grained si basa sui claim del JWT:
\texttt{preferred\_username}, \texttt{Ruolo}, \texttt{Direzione} sono
propagati dai sistemi HR a Entra ID, e la lista dei gruppi (campo
\texttt{group Vera AI}) è quanto consente l'accesso selettivo ai dati
sensibili. CooPolicy in particolare confronta i gruppi del chiamante con
una \emph{whitelist} configurata via env (\texttt{DOCS\_ALLOWED\_GROUPS})
prima di accettare la richiesta: se l'utente non vi appartiene, riceve
un 403 senza che l'esistenza dei corpus venga rivelata.

\subsection{Catena di sicurezza pre e post generazione}
\label{subsec:agenti-security}

Coerentemente con il principio di governance facilitata, ogni invocazione di
modello attraversa una catena di controlli predisposta nel servizio
\texttt{model\_armor\_service.py} di Vera, e replicata in forma più snella
negli altri agenti. La catena ha tre stadi:

\begin{enumerate}
  \item \textbf{Sanitizzazione del prompt.} Prima di inviare il messaggio
        all'LLM, il testo passa per il template di Google Cloud Model
        Armor~\cite{gcp_modelarmor}, che combina un classificatore di
        \emph{Responsible AI} (categorie \emph{hate}, \emph{harassment},
        \emph{sexually explicit}, \emph{dangerous}) con il rilevatore
        di \emph{prompt injection}~\cite{perez2022promptinjection} e con
        la \emph{Sensitive Data Protection} di Google Cloud~\cite{gcp_dlp}
        per le PII. I tipi PII rilevanti per il contesto italiano (Codice
        Fiscale, Partita IVA, IBAN) sono configurati come \emph{custom
        info types}; quando la PII viene rilevata, il backend non
        sostituisce ciecamente il valore ma rifiuta la richiesta con un
        messaggio che invita l'utente a riformulare. Il rifiuto è
        appropriato perché l'autenticità del dato originale è spesso
        necessaria per l'utilità della risposta, e il \emph{tokenizing}
        cieco produrrebbe risposte inutilizzabili.
  \item \textbf{Esenzioni per gruppo.} Alcuni gruppi (ad esempio l'ufficio
        legale) hanno legittimo bisogno di trattare PII specifiche.
        Una tabella PostgreSQL \texttt{model\_armor\_group\_exemptions}
        e una \texttt{pii\_group\_exemptions} consentono di escludere
        per gruppo singoli filtri (ad esempio \texttt{prompt\_injection}
        o singole categorie RAI) o singoli tipi di PII. Al momento
        della verifica, il backend interseca la lista di esenzioni
        attive per l'utente con i \emph{reasons} restituiti da Model
        Armor: se ogni motivo è coperto da esenzione, la richiesta
        prosegue.
  \item \textbf{Sanitizzazione della risposta e audit.} La stessa funzione
        di sanitizzazione viene applicata all'output del modello prima del
        salvataggio in conversazione, e ogni interazione è registrata in
        modo strutturato nella tabella \texttt{chat\_logs} con
        \emph{action}, modello selezionato, conteggio token in/out e flag
        di citazioni. Eventi di blocco vengono inoltre registrati su un
        \emph{security log} dedicato (\texttt{security\_logger.log\_prompt\_blocked},
        \texttt{log\_response\_blocked}), che alimenta la dashboard di
        amministrazione descritta in Sez.~\ref{sec:agenti-admin-impl}.
\end{enumerate}

Il consumo di token è controllato dal servizio di quota: il backend stima
preventivamente il consumo della richiesta, lo confronta con il budget mensile
del badge dell'utente (\emph{Bronze}, \emph{Silver}, \emph{Gold}) e blocca
con un 400 se il limite sarebbe superato. Al termine della chiamata il
consumo effettivo viene contabilizzato. Questo è il meccanismo concreto che
realizza il principio di \emph{Implementazione sostenibile} del
Cap.~\ref{chap:agenti-architettura}.

% =====================================================================
\section{Vera --- backend e frontend}
\label{sec:agenti-vera-impl}
% =====================================================================

Vera è il backend più articolato; le scelte realizzative significative
ruotano attorno a quattro problemi: \emph{routing} del modello, \emph{streaming}
SSE robusto, \emph{Web grounding} con citazioni inline, gestione dei documenti
caricati dall'utente.

\subsection{Routing del modello: badge e complessità della query}
\label{subsec:vera-routing}

La selezione del modello è governata dal modulo \texttt{model\_router.py}:valuta
euristicamente la complessità della richiesta sommando tre indicatori
(lunghezza superiore a 500~caratteri, presenza di parole-chiave riconducibili
ad analisi articolate, presenza di più punti interrogativi) e instrada verso
Pro se almeno due indicatori sono positivi. Se nel database è configurato un
mapping esplicito badge--modello (tabella \texttt{badge\_model\_access}), la
logica DB ha precedenza sull'euristica, in modo che gli amministratori
possano modificare il comportamento senza un nuovo deploy. La scelta di
mantenere il \emph{fallback} euristico anche in presenza della tabella DB ha
un costo di poche righe di codice e mette al riparo il sistema da
configurazioni vuote in ambienti freschi.

Questa logica realizza in forma essenziale, e per uso interno al singolo
agente, il \emph{cost-quality trade-off} formalizzato dalla letteratura sul
\emph{LLM routing}~\cite{chen2023frugalgpt,ong2024routellm}: un modello più
piccolo è preferito quando è sufficiente, un modello più grande è impiegato
solo quando la query lo giustifica.

\subsection{Streaming SSE con sessioni dedicate e keepalive ticker}
\label{subsec:vera-sse}

L'endpoint \texttt{/chat/stream} restituisce una \texttt{StreamingResponse}
FastAPI di tipo \texttt{text/event-stream}. La struttura del generatore
asincrono che la alimenta è il risultato di un'iterazione architetturale
non banale: la prima versione, che condivideva la \texttt{AsyncSession}
SQLAlchemy fra il pre-\emph{check} di quota e il loop di streaming,
provocava \emph{rollback} spuri al secondo messaggio della stessa
conversazione, perché FastAPI chiudeva la sessione iniettata al ritorno
del \texttt{StreamingResponse} mentre il generatore era ancora attivo.
La soluzione adottata è separare due sessioni: \texttt{pre\_db}, di vita
breve, esegue tutti i controlli sincroni (quota, sanitizzazione del
prompt via Model Armor, scelta del modello, \emph{lookup} della
conversazione, caricamento della history) e viene chiusa subito; una
seconda sessione \texttt{stream\_db} è creata all'interno del
generatore e vive per tutta la durata dello stream. I dati che
attraversano il confine tra le due sessioni sono primitivi (UUID,
stringhe, dizionari), così da non innescare \emph{lazy loading} con
sessione chiusa.

Un secondo accorgimento è il \emph{keepalive ticker}: ogni 3~secondi un
task asincrono pubblica nella coda di eventi un commento SSE
(\texttt{": keepalive\textbackslash n\textbackslash n"}) che attraversa il
flusso senza essere rappresentato lato client. Senza questo dispositivo,
proxy e browser (in particolare Firefox) chiudono la connessione durante
le pause prolungate del modello (primo token, retrieval, rate limit
transienti). La coda di eventi è una \texttt{asyncio.Queue} che multiplex
flussi di dati e ticker; un sentinella \texttt{None} segnala la fine
dello stream e il task di \emph{ticker} viene cancellato in finally.

\subsection{Pattern di esecuzione: documenti vs Web grounding}
\label{subsec:vera-paths}

Il generatore distingue due percorsi mutuamente esclusivi sulla base della
presenza di documenti caricati nella conversazione. Quando ci sono \emph{upload}
in stato \texttt{ready}, viene costruita o riusata una \emph{Gemini Context
Cache} associata alla conversazione, che evita di ricaricare i medesimi
documenti a ogni turno; la generazione avviene quindi tramite l'agente ADK
\texttt{doc\_qa\_agent}, configurato in modo da utilizzare la cache come
contenuto pre-cacheato (\texttt{generate\_content\_config.cached\_content}).
In assenza di upload, il generatore percorre il \emph{path} di conversazione
generale, che monta un \texttt{LlmAgent} ADK con
l'\texttt{EnterpriseWebSearchTool} integrato: il modello decide
autonomamente, in base alla domanda, se invocare la \emph{search} o
rispondere con la sola conoscenza parametrica.

In entrambi i percorsi la frammentazione dell'output è gestita con la modalità
\texttt{StreamingMode.SSE} di ADK: il \emph{loop} di eventi del Runner emette
\texttt{Event} parziali, di cui vengono propagate solo le \texttt{parts} con
campo \texttt{partial = True}, escludendo i blocchi marcati come
\texttt{thought}. Quando il modello restituisce \texttt{finish\_reason}
uguale a \texttt{MAX\_TOKENS}, il backend invia automaticamente un turno
addizionale ``Continua la risposta dal punto esatto in cui si è interrotta''
fino a una soglia di sicurezza (\texttt{\_MAX\_CONTINUATIONS = 4}), così
da rendere la risposta completa anche quando supera il limite di output del
modello.

\subsection{Citazioni inline e \emph{grounding metadata}}
\label{subsec:vera-citations}

Quando il modello invoca la \emph{Web search}, l'evento ADK porta in un campo
dedicato (\texttt{event.grounding\_metadata}) la lista dei chunk consultati,
con URL, titolo e \emph{snippet}. La funzione \texttt{extract\_web\_citations}
trasforma questa struttura in un array di citazioni che il frontend rende
inline, cliccabili dall'utente. Per i documenti caricati la citazione viene
estratta da un \texttt{citation\_extractor} che attraversa il testo finale e
risolve i riferimenti contro la lista dei \texttt{cache\_uploads}. Il fatto
che la fonte sia sempre verificabile e cliccabile, e che la sua presenza sia
imposta dal sistema (non opzionale per il modello), realizza il vincolo di
verificabilità descritto in
Cap.~\ref{chap:agenti-architettura}, Sez.~\ref{subsec:vera-caso-uso}.

% =====================================================================
\section{CooPolicy --- backend}
\label{sec:agenti-coopolicy-impl}
% =====================================================================

CooPolicy è la realizzazione ``RAG-native'' della suite: tutto il backend
ruota attorno alla pipeline di Retrieval-Augmented
Generation~\cite{lewis2020rag}, con scelte mirate a comprimere il rumore di
recupero e ad ancorare la generazione esclusivamente al materiale recuperato.

\subsection{Multi-corpus e classificazione della query}
\label{subsec:coopolicy-multicorpus}

Una prima scelta architetturalmente significativa è la rinuncia al
\emph{single corpus}. La documentazione interna è ripartita in dieci corpus
Vertex AI~\cite{gcp_vertex_rag} tematici (\texttt{policy-hr},
\texttt{contratti-smartworking}, \texttt{amministrazione-personale},
\texttt{procedure-aziendali}, \texttt{processi-regolamenti},
\texttt{governance-identita}, \texttt{sicurezza},
\texttt{punti-vendita-operativita}, \texttt{materiale-pdv},
\texttt{comunicazioni-generali}), ciascuno con un proprio descrittore in
italiano del dominio coperto. La motivazione è duplice: da un lato la
qualità del retrieval cala rapidamente al crescere della cardinalità dei
chunk indicizzati nello stesso spazio vettoriale; dall'altro, alcune
classi di documenti hanno regole di interpretazione
specifiche (i contratti collettivi vincolano la risposta in modo diverso
dalle policy aziendali) e mantenerle separate facilita l'evoluzione
indipendente.

A monte del retrieval, una funzione \texttt{classify\_query} chiede a
Gemini~2.5 Flash di scegliere il \emph{corpus primario} ed eventualmente
un \emph{corpus secondario} fra quelli configurati, restituendo
strettamente un oggetto JSON con due chiavi. La scelta di Gemini Flash è
motivata dalla coppia di richieste essenziale di questo passo: bassa
latenza (sub-secondo) e basso costo unitario, condizione necessaria
perché il classificatore venga invocato a ogni richiesta senza pesare
sull'economia del sistema. Il prompt include il descrittore di ciascun
corpus disponibile.

\subsection{Retrieval ibrido e soglia di rilevanza}
\label{subsec:coopolicy-retrieval}

Il retrieval è eseguito in parallelo sui corpus selezionati attraverso
\texttt{rag.retrieval\_query} di Vertex AI RAG Engine~\cite{gcp_vertex_rag},
con \texttt{top\_k = 40} sul corpus primario e \texttt{top\_k = 20} sul
secondario quando presente. La modalità di ricerca preferita è
\emph{hybrid} (\texttt{rag.HybridSearch} con \texttt{alpha = 0.5}), che
bilancia ricerca vettoriale densa e matching lessicale BM25:
l'\emph{alpha} 0.5 è stato scelto per bilanciare due esigenze
spesso in tensione, la robustezza semantica della ricerca densa e la
precisione su nomi propri, codici di policy e termini tecnici esatti
(es. ``CCNL'', ``RGPD'', codici di reparto), per i quali il segnale BM25
è insostituibile. Quando il corpus non supporta la modalità ibrida, il
codice ricade in modo automatico sulla sola ricerca vettoriale, evitando
errori di parametri.

I chunk restituiti sono filtrati con una soglia
\texttt{RAG\_SCORE\_THRESHOLD~=~0.15} sulla distanza coseno: il valore è
deliberatamente permissivo, in coerenza con la decisione di delegare al
modello generativo (vincolato dal \emph{system prompt}) la decisione su
quali chunk siano effettivamente utili. L'aspettativa è che un modello in
modalità a bassa temperatura (Sez.~\ref{subsec:coopolicy-generation})
sa rifiutare contesto irrilevante, mentre una soglia troppo restrittiva
rischia di escludere materiale marginalmente rilevante che però contiene
la risposta. I risultati di primario e secondario vengono uniti,
deduplicati per testo identico, ordinati per distanza coseno crescente
e troncati a \texttt{RAG\_TOP\_K} chunk totali.

Una linea di sviluppo recente è la richiesta di \emph{query expansion}
ispirata a~\cite{ma2023query_expansion} e a HyDE~\cite{gao2022hyde}. Nella
versione attuale di CooPolicy non è ancora attiva: la combinazione
multi-corpus + hybrid + soglia bassa ha mostrato sufficiente \emph{recall}
per il dominio in esame; restano disponibili come direttrice di
miglioramento in caso si registri un calo di \emph{recall} con la
crescita del corpus.

\subsection{Arricchimento dei chunk e gestione dei metadati}
\label{subsec:coopolicy-meta}

Ogni PDF nel bucket GCS è accompagnato da un sidecar \texttt{.meta.json}
generato in fase di ingestione, contenente titolo leggibile e data del
documento. La funzione \texttt{\_enrich\_chunks\_with\_meta} scarica in
parallelo i sidecar dei chunk recuperati, con cache in-memory a TTL
un'ora, e arricchisce ciascun chunk con il titolo umano-leggibile e con la
data; quest'ultima è poi utilizzata dal modello per privilegiare la
fonte più recente in caso di ridondanza informativa. I file di metadati
(\texttt{.json}) e i tabulari grezzi (\texttt{.csv}), che possono
finire indicizzati per altri motivi, sono filtrati esplicitamente prima
di costruire il prompt, evitando che il modello citi rumore.

\subsection{Generazione vincolata e citazioni}
\label{subsec:coopolicy-generation}

Il prompt di sistema impone, in italiano, sette regole non negoziabili:
analisi integrale dei documenti recuperati prima di rispondere; risposta
basata esclusivamente sul materiale fornito; obbligatoriet\`a delle
citazioni \texttt{[N]} inline; massima precisione su numeri, articoli e
condizioni; divieto esplicito di parafrasi libere; tono professionale
italiano; consapevolezza temporale (preferire la fonte più recente,
avvisare se il documento ha più di due anni). Il modello è Gemini~2.5 Pro
con \texttt{temperature = 0.1} e \texttt{max\_output\_tokens = 8192}: la
temperatura quasi zero massimizza l'aderenza al prompt e riduce
ulteriormente il rischio di allucinazione~\cite{xu2024hallucination}, in
linea con la funzione di un agente di tipo ``ricerca verificabile''.

Il post-processing della risposta è il dispositivo che garantisce
l'integrità delle citazioni. Le citazioni multiple
(\texttt{[1, 2, 3]}) vengono espanse in citazioni singole
(\texttt{[1][2][3]}); le citazioni con indice fuori dal range dei
documenti realmente forniti --- ovvero le citazioni allucinate ---
vengono rimosse dal testo finale; gli indici sono infine rinumerati
in ordine di prima comparsa, in modo che la lista delle fonti
restituita al frontend rispecchi l'ordine di lettura. Per ciascuna
citazione viene anche estratto uno \emph{snippet} di 300~caratteri
centrato sulle keyword della query, utile per la \emph{preview} nel
frontend.

\subsection{ACL e proxy autenticato per i documenti}
\label{subsec:coopolicy-acl}

Il primo passo di ogni endpoint sensibile è
\texttt{\_require\_docs\_group}: l'utente deve appartenere ad almeno un
gruppo Entra ID inserito nella \emph{whitelist}
\texttt{DOCS\_ALLOWED\_GROUPS}, altrimenti la richiesta è respinta con
403 senza ulteriori dettagli. L'accesso al PDF originale, infine, non
avviene tramite URL pubblico: l'endpoint \texttt{/api/docs/file} agisce
come proxy autenticato che, verificato il token, ottiene un \emph{access
token} cloud-platform tramite le credenziali del service account del
servizio e scarica il binario direttamente dall'API di GCS, restituendolo
con \texttt{Content-Disposition: inline}. Il modello non ha mai
visibilit\`a degli URL firmati e l'utente non pu\`o aggirare la
verifica del token visitando un link.

% =====================================================================
\section{AllyCare --- backend}
\label{sec:agenti-allycare-impl}
% =====================================================================

AllyCare è la realizzazione del paradigma Text-to-SQL su BigQuery per
l'analisi dei sondaggi NPS. La traduzione di linguaggio naturale in SQL
è un campo di ricerca attivo, con benchmark
consolidati~\cite{yu2018spider,li2023bird} e tecniche di decomposizione e
auto-correzione che migliorano l'accuratezza su query
complesse~\cite{pourreza2023dinsql}. L'implementazione di AllyCare
combina questi spunti con il paradigma agentico
ReAct~\cite{yao2022react}, esponendo al modello tre \emph{tool} di
interazione con i dati e affidando all'\emph{agentic loop} di ADK la
gestione del ciclo Reason--Act--Observe.

\subsection{Iniezione selettiva dello schema}
\label{subsec:allycare-schema}

Il prompt di sistema include lo schema BigQuery delle due tabelle
rilevanti (\texttt{mapp\_nps\_survey\_internal\_prod} con i dati di
sondaggio grezzi e \texttt{mapp\_nps\_survey\_ai\_enriched\_prod} con i
campi prodotti da una pipeline di arricchimento AI a monte) ottenuto via
\texttt{bq\_client.get\_table} e mantenuto in cache per un'ora. Allo
schema tecnico viene affiancato un \emph{glossario semantico} in
italiano: per ciascuna colonna è descritto il significato funzionale
(es. \texttt{voto\_item01: ``Il personale del negozio è gentile e
competente''; valori -2 = mancata risposta, -1 = non lo so, 0--9 = voto}).
Insieme allo schema viene iniettata anche la lista esatta delle
descrizioni di negozio (\texttt{store\_desc}, in cache 24~ore) tratta
direttamente da BigQuery, in modo che il modello possa risolvere
varianti colloquiali (``Sasso'', ``Iper Bologna'') in match esatti senza
allucinare codici. Questa strategia di \emph{schema injection
selettiva}, in luogo dell'esposizione del catalogo completo del data
warehouse, è coerente con la letteratura
Text-to-SQL~\cite{pourreza2023dinsql}: meno colonne il modello vede,
meno opportunità ha di confondersi.

Il prompt impone inoltre regole di consistenza SQL specifiche del
dominio: i filtri sui voto-item devono includere \texttt{voto\_itemXX
>= 0} nella \texttt{WHERE} (escludendo i marcatori \texttt{-1} e
\texttt{-2}); tutti i confronti su campi categorici testuali devono
usare \texttt{UPPER(TRIM(...))} per neutralizzare incoerenze di
\emph{case}; la distinzione fra \texttt{LIKE} e uguaglianza esatta
sui nomi di negozio segue regole esplicite. L'effetto è che AllyCare
tipicamente non sbaglia i filtri ``stupidi'' che dominano gli errori in
Text-to-SQL applicato in dominio reale.

\subsection{Tre \emph{tool} e l'\emph{agentic loop} ADK}
\label{subsec:allycare-tools}

Lo strato di tool è dichiarato in \texttt{tools.py} ed esporta tre
funzioni Python registrate come \emph{ADK tools}:

\begin{description}
  \item[\texttt{run\_sql(query)}]
    Esegue una query SQL su BigQuery con un \emph{guard} sintattico: la
    query è parsificata con \texttt{sqlparse} e respinta se contiene
    \texttt{DDL} o \texttt{DML} di scrittura
    (\texttt{INSERT}/\texttt{UPDATE}/\texttt{DELETE}); se non già
    presente, viene aggiunto un \texttt{LIMIT~5000}; il job è eseguito
    con \texttt{maximum\_bytes\_billed = 10\textsuperscript{9}} per
    bloccare query accidentalmente costose. Il risultato è restituito
    al modello come tabella Markdown delle prime 750~righe, con
    indicazione del totale; la copia completa del DataFrame è
    memorizzata nello \texttt{tool\_context.state} per essere
    riutilizzata da \texttt{generate\_chart} senza una seconda
    interrogazione. Se la query produce un errore SQL, il messaggio
    viene impacchettato come \texttt{function\_response} e il modello
    può auto-correggersi al ciclo successivo: questo è il meccanismo
    concreto di \emph{self-correction}~\cite{pourreza2023dinsql}.
  \item[\texttt{analyze\_comments(query, question)}]
    Estrae un campione di commenti testuali grezzi mediante una query
    dedicata, ne preleva un sotto-campione casuale di al più 200 record
    e chiede a Gemini~2.5 Pro una sintesi qualitativa con citazioni
    \emph{verbatim}. Il prompt di analisi è scritto in modo da essere
    integrato nella risposta più ampia e a non duplicare l'intestazione
    di sezione. È usato esclusivamente per analisi qualitative: per
    qualunque dato quantitativo (sentiment, conteggi di temi, percentuali)
    il prompt di sistema impone \texttt{run\_sql} sulla tabella
    \texttt{ai\_enriched\_prod}.
  \item[\texttt{generate\_chart(visualization\_goal)}]
    Costruisce un grafico Matplotlib/Seaborn a partire dal DataFrame
    cachato dall'ultimo \texttt{run\_sql}, scegliendo la rappresentazione
    in funzione di parole-chiave nel \emph{goal}
    (\texttt{torta}/\texttt{pie} $\rightarrow$ pie chart;
    \texttt{linee}/\texttt{trend} $\rightarrow$ line chart; default
    $\rightarrow$ bar chart). Il grafico è codificato in PNG, convertito
    in base64 con prefisso \emph{data URI} e iniettato nel session state.
    Il backend lo restituisce poi nella risposta come campo
    \texttt{chart\_and\_data}, lasciando il rendering alla logica del
    frontend. Il chart viene incluso nel payload finale solo se
    \texttt{generate\_chart} è stato effettivamente invocato nel turno,
    evitando di restituire il grafico del turno precedente.
\end{description}

L'\emph{agentic loop} è governato dal \texttt{Runner} ADK che, a partire
dall'\texttt{LlmAgent} configurato con i tre tool e con il prompt di
sistema, alterna in modo trasparente fasi di ragionamento e chiamate a
tool fino a quando il modello restituisce la risposta finale. Il
modello scelto per il turno (Pro o Flash) è iniettato nell'agente al
momento della costruzione, e il prompt di sistema include il nome
proprio dell'utente e la sua direzione (estratti dal token Entra ID),
così che il tono della risposta sia personalizzato.

Il modulo \texttt{select\_model} replica la logica già descritta per Vera
in forma adattata al dominio: parole-chiave segnaletiche di analisi
articolate (\texttt{commenti}, \texttt{sentiment}, \texttt{trend},
\texttt{breakdown}, \texttt{voto\_item}) instradano verso Gemini~2.5
Pro; richieste meramente quantitative (\texttt{quanti rispondenti},
\texttt{media}, \texttt{nps medio}) verso Flash. Il default conservativo
è Pro, perché un calo di qualità su domanda complessa è più visibile
all'utente di un \emph{overspending} marginale.

\subsection{Disciplina del ragionamento e garanzie anti-allucinazione}
\label{subsec:allycare-discipline}

Il prompt di sistema impone una disciplina di output esplicita: ogni
risposta deve aprirsi con una sezione \texttt{\#\#\# RAGIONAMENTO} che
enuncia obiettivo, scelta delle colonne, strategia SQL e nota
anti-allucinazione; deve quindi chiamare il \emph{tool} appropriato
prima di pronunciare numeri; deve infine produrre la sezione
\texttt{\#\#\# RISPOSTA} solo dopo aver ricevuto i risultati. La
disciplina è una forma di Chain-of-Thought esplicito~\cite{wei2022cot},
estesa con un protocollo a due tag che separano in modo riconoscibile la
fase di pianificazione dall'output utente. Una funzione
\texttt{\_clean\_answer} elimina il blocco di ragionamento da quanto
viene effettivamente mostrato all'utente, in modo che la
trasparenza sia disponibile in audit (le righe del \emph{thought}
restano nei \emph{log} ADK) ma l'esperienza utente sia pulita.

La ``gerarchia della verità'' codificata nel prompt impone che i dati
del database siano la sola fonte autoritativa: il modello non può
classificare commenti per produrre numeri (la categorizzazione
quantitativa va sempre dalla tabella arricchita). Una regola di
``fiducia nel tool'' impone inoltre che, se \texttt{run\_sql} ha
restituito un dato, sia quel dato a essere riportato, anche se il
modello ricorda numeri diversi da turni precedenti. Questo accorgimento
ha eliminato in fase di test una classe di errori in cui il modello,
ricalcolando mentalmente, contraddiceva il proprio tool.

\subsection{Resilienza operativa}
\label{subsec:allycare-resilience}

Il loop ADK è racchiuso in un \texttt{try/except} che riconosce gli
errori transitori della piattaforma di inferenza
(\texttt{RESOURCE\_EXHAUSTED}, \texttt{429}, \texttt{503}, \emph{rate
limit}) e li ripete con \emph{backoff} esponenziale fino a tre
tentativi, evitando che picchi di traffico contro Vertex AI degradino
la qualità di servizio percepita dall'utente. La medesima strategia di
retry transitorio è utilizzata anche da CooPolicy nelle chiamate al RAG
Engine. Le query catturate dai tool sono accumulate nello
\texttt{state} ADK e restituite al frontend come campo
\texttt{queries}: l'utente può così ispezionare le SQL effettivamente
eseguite, abilitando un meccanismo di trasparenza che è anche un
deterrente di fatto contro errori silenziosi.

% =====================================================================
\section{SysAid --- backend Gemini e backend Sonnet}
\label{sec:agenti-sysaid-impl}
% =====================================================================

SysAid è l'implementazione su cui è stato condotto il primo esperimento
controllato di indipendenza dal fornitore di LLM, costruendo due backend
funzionalmente equivalenti, alimentati rispettivamente da Gemini~2.5
Pro/Flash su Vertex AI e da Claude Sonnet~4.6 per il tramite di
\texttt{AnthropicVertex}. La scelta di mantenere distinti due backend,
invece di astrarre un unico backend con \emph{driver} polimorfi sul
modello, è stata dettata dall'obiettivo dell'esperimento: rendere
visibili in produzione, e per gli stessi utenti finali, le differenze
di comportamento dei due modelli, evitando che lo strato di astrazione
ne mascherasse le idiosincrasie. La presenza di entrambi prepara inoltre
il terreno alla valutazione comparativa di Orion
(Cap.~\ref{chap:orion-valutazione}), in linea con la metodologia
\emph{LLM-as-a-Judge}~\cite{zheng2023llmjudge}.

\subsection{Architettura comparativa}
\label{subsec:sysaid-arch-comp}

L'architettura comparativa minimizza l'attrito sperimentale: schema
BigQuery condiviso (vista \texttt{service\_req\_view} sulla base ticket
SysAid), interfaccia React identica, prompt di sistema identico nei
contenuti, set di tool funzionalmente equivalente (\texttt{run\_sql},
\texttt{analyze\_text}, \texttt{generate\_chart}), conversazioni
persistite su tabelle separate per non confondere i due flussi.

\begin{table}[htbp]
  \centering
  \caption[Confronto implementativo SysAid Gemini vs.\ SysAid Sonnet]{Confronto
    implementativo tra le due varianti di SysAid. Le righe evidenziate
    mostrano le differenze tecniche significative.}
  \label{tab:sysaid-confronto}
  \begin{tabular}{lll}
    \toprule
    \textbf{Dimensione}      & \textbf{SysAid Gemini}        & \textbf{SysAid Sonnet}            \\
    \midrule
    Modello LLM              & Gemini 2.5 Flash / Pro        & Claude Sonnet 4.6                 \\
    Piattaforma              & Vertex AI (\texttt{europe-west8}) & AnthropicVertex (\texttt{europe-west1}) \\
    SDK                      & \texttt{vertexai.generative\_models} & \texttt{anthropic.AnthropicVertex} \\
    \textbf{Schema dei tool} & \textbf{\texttt{FunctionDeclaration}} & \textbf{JSON Schema}     \\
    \textbf{Stop reason tool} & \textbf{\texttt{function\_call} part} & \textbf{\texttt{stop\_reason = tool\_use}} \\
    \textbf{Adattamento}     & \textbf{Nessuno}              & \textbf{Thin adapter dedicato}    \\
    Schema BigQuery          & Condiviso                     & Condiviso                          \\
    Prompt di sistema        & Identico                      & Identico                           \\
    Tool disponibili         & \texttt{run\_sql}, \texttt{analyze\_text}, \texttt{generate\_chart}
                                                             & \texttt{run\_sql}, \texttt{analyze\_text}, \texttt{generate\_chart} \\
    Frontend React           & Condiviso                     & Condiviso                          \\
    \bottomrule
  \end{tabular}
\end{table}

Una nota architetturale non banale è la differenza di regione GCP fra le
due varianti: Gemini è disponibile in \texttt{europe-west8}, mentre
Claude Sonnet su Vertex è offerto al momento solo in
\texttt{europe-west1}. Entrambe le regioni si trovano nel perimetro
europeo e soddisfano il requisito di sovranità del dato.

\subsection{Variante Gemini: \emph{function calling} nativo}
\label{subsec:sysaid-gemini}

La variante \texttt{sysaid-backend} dichiara i tre tool con
\texttt{FunctionDeclaration} di \texttt{vertexai.generative\_models},
li raggruppa in un oggetto \texttt{Tool} e li passa al
\texttt{GenerativeModel} insieme a \texttt{GENERATION\_CONFIG = \{
"temperature": 0 \}}. Il \emph{loop} di ragionamento è scritto a mano:
ad ogni risposta del modello, il backend ispeziona le \texttt{parts} alla
ricerca di \texttt{function\_call}, esegue la funzione corrispondente,
costruisce un \texttt{Part.from\_function\_response} con il risultato e
lo aggiunge all'history per il turno successivo, fino a un cap di
\texttt{MAX\_REASONING\_STEPS = 20} iterazioni. La temperatura zero,
in continuità con CooPolicy e con la letteratura sulla
calibrazione~\cite{kadavath2022knowsknows}, riduce la varianza
dell'output e rende riproducibile il comportamento del modello su query
identiche.

\subsection{Variante Sonnet: \emph{thin adapter} per Anthropic}
\label{subsec:sysaid-sonnet}

La variante \texttt{sysaid-sonnet-backend} accede a Claude Sonnet
attraverso \texttt{AnthropicVertex}, l'integrazione che permette di
chiamare i modelli Anthropic mantenendo l'autenticazione e la fatturazione
all'interno del progetto GCP~\cite{anthropic_docs}. L'API Anthropic
attende una semantica diversa per i tool e per l'history dei messaggi
rispetto a Vertex AI nativo, e questo richiede l'introduzione di un
\emph{thin adapter} che si fa carico delle differenze.

\paragraph{Definizione dei tool.} Mentre Gemini accetta
\texttt{FunctionDeclaration} con \texttt{parameters} in formato
proprietario, Anthropic richiede una lista di oggetti con campi
\texttt{name}, \texttt{description} e \texttt{input\_schema} in
JSON~Schema standard. La lista \texttt{SYSAID\_TOOLS} è quindi
duplicata in JSON~Schema; il sotto-insieme \texttt{SYSAID\_TOOLS\_REPORT}
omette \texttt{generate\_chart} ed è impiegato dall'endpoint di report
giornalieri (cfr. canale WhatsApp), che non può rendere immagini.

\paragraph{Loop a tool e \texttt{stop\_reason}.} Anthropic comunica
l'invocazione di un tool interrompendo la generazione con
\texttt{stop\_reason = "tool\_use"}. La risposta contiene blocchi
tipizzati, fra cui zero o più \texttt{tool\_use} con \texttt{id},
\texttt{name} e \texttt{input}. Il loop ripete fino a venti iterazioni:
per ciascun blocco \texttt{tool\_use} esegue la funzione corrispondente,
serializza il risultato in un blocco \texttt{tool\_result} con lo stesso
\texttt{tool\_use\_id} e lo accoda come messaggio \emph{user}, prima di
richiamare il modello. Quando \texttt{stop\_reason} è \texttt{end\_turn}
e nessun tool è stato invocato in tutto il turno, un \emph{guardrail}
specifico verifica che la sezione \texttt{\#\#\# RISPOSTA} non sia
presente: se lo è, il modello ha tentato di rispondere senza chiamare
\texttt{run\_sql}, e il backend forza un retry con un messaggio di
correzione. È il riflesso anti-allucinazione del prompt, applicato lato
controller invece che lato modello.

\paragraph{Serializzazione dei content block.} L'SDK Anthropic
restituisce content block che includono campi interni
(\texttt{caller}, \texttt{parsed\_output}, \ldots) non accettati dall'API
quando il blocco viene rispedito come parte dell'history. Per evitare
errori 400, l'helper \texttt{\_content\_to\_api} applica una
\emph{whitelist} per tipo di blocco
(\texttt{text}, \texttt{tool\_use}, \texttt{tool\_result},
\texttt{image}) che mantiene solo i campi documentati. Senza questa
\emph{whitelist} le iterazioni successive del loop fallirebbero in modo
ricorrente.

\paragraph{Gestione di output lunghi e \emph{streaming} interno.} Per
sostenere \texttt{max\_tokens = 32000} senza incorrere in timeout
dell'SDK, il backend usa internamente l'API \emph{streaming} di
Anthropic (\texttt{messages.stream}), accumulando il messaggio finale
con \texttt{stream.get\_final\_message()} prima di processarlo nel
loop. Quando \texttt{stop\_reason = "max\_tokens"}, il backend aggiunge
un messaggio utente ``Continua\ldots'' e ripete il giro, in
analogia con il meccanismo di continuation di Vera
(Sez.~\ref{subsec:vera-paths}). Se il loop esce per esaurimento del
budget di iterazioni con \texttt{\#\#\# RISPOSTA} ancora assente, un
\emph{turno di sintesi forzata} chiede al modello di produrla con i
dati gi\`a raccolti, evitando finestre vuote per l'utente.

L'effetto netto del \emph{thin adapter} è che la logica di business
condivisa (esecuzione SQL, analisi testuale, grafici, persistenza)
risiede in una sola parte del backend e non viene duplicata; gli
adattamenti rispetto al fornitore restano localizzati alla traduzione
di schema dei tool, all'ispezione dei \texttt{stop\_reason} e alla
serializzazione dei content block. La struttura ha permesso di
rispettare il principio di \emph{Implementazione agnostica e
sostenibile} senza aprire una refactoring globale degli agenti già in
produzione.

% =====================================================================
\section{Pannello di amministrazione}
\label{sec:agenti-admin-impl}
% =====================================================================

Il monitoraggio e la configurazione dell'intera suite sono affidati a un
servizio dedicato (\texttt{admin-backend} + \texttt{admin-frontend})
deliberatamente isolato dagli applicativi rivolti all'utente finale. Il
backend FastAPI espone i propri router sotto il prefisso
\texttt{/api/admin/} e \`e composto da moduli verticali ciascuno
focalizzato su un dominio: gestione utenti, quote per badge, statistiche
di consumo, log audit, feedback strutturato, mapping
badge--modelli, piani di rate limit, revisioni, conversation starter,
impostazioni globali (es. temperatura), gruppi di esenzione PII e Model
Armor. Il frontend React effettua un proprio login MSAL e, sul
\emph{load balancer}, è instradato sul sub-path \texttt{/admin/}
attraverso una regola di URL map dedicata, in modo da non condividere
sessione e bundle JavaScript con il portale. Questa segregazione
contiene il \emph{blast radius} di eventuali vulnerabilit\`a del
pannello.

Le funzioni di osservabilità sono già pensate in vista dell'evoluzione
con Orion. La tabella \texttt{chat\_logs} (popolata da Vera, replicata
con piccole variazioni dagli altri agenti) e i log strutturati di
\texttt{security\_logger} costituiscono il fondamento dello stesso
\emph{audit trail} che, una volta arricchito con il record di
\texttt{orion\_routing\_logs} (Cap.~\ref{chap:orion-implementazione}),
abiliterà la dashboard \emph{RouterMetrics} per la valutazione
sperimentale della Sez.~\ref{chap:orion-valutazione}. Nessun nuovo
servizio sarà perciò necessario per ospitare la telemetria
dell'orchestratore.
